
\documentclass[manuscript]{acmart}
\setcopyright{none}
\settopmatter{printacmref=false,printfolios=false}
\acmJournal{CSUR}
\acmYear{2026}
\title{ACM Family: Official ACM Submission Preview}
\author{A. Placeholder}
\email{verified@example.org}
\affiliation{%
  \institution{Department of Verified Templates, Placeholder Institute}
  \city{Placeholder City}
  \country{Country}
}
\author{B. Example}
\affiliation{%
  \institution{Journal Systems Lab, Example University}
  \city{Example City}
  \country{Country}
}
\begin{document}
\begin{abstract}
This submission-style preview is rendered with the official ACM acmart template in manuscript mode. The content remains placeholder text, but the page is intentionally organized as a realistic review-format article so the ACM-specific front matter, section hierarchy, figure flow, table rhythm, CCS metadata, and bibliography style are easier to inspect.
\end{abstract}
\keywords{ACM template, manuscript review format, official preview, figure placement, table placement}
\begin{CCSXML}
<ccs2012>
<concept>
<concept_id>10010147.10010178</concept_id>
<concept_desc>Computing methodologies</concept_desc>
<concept_significance>500</concept_significance>
</concept>
</ccs2012>
\end{CCSXML}
\ccsdesc[500]{Computing methodologies}
\maketitle
\section{Introduction}
The introduction is written as compact technical prose so the manuscript preview reads like a plausible ACM journal submission rather than a bare template check page. This makes the relation between title block, abstract, CCS concepts, and early body text easier to judge.
\section{Related Work}
The related-work section remains explicit because it helps reveal how the ACM manuscript layout handles dense section transitions, literature framing, and survey-style or method-style context.
\section{Methodology}
The methodology section introduces representative mathematics and a figure placeholder so the acmart review-format page can be inspected under more realistic technical content.
\begin{equation}
L = \sum_{i=1}^{M} w_i \left\| y_i - \hat{y}_i \right\|_2^2.
\end{equation}
\begin{figure}[t]
\centering
\fbox{\parbox[c][2.0in][c]{0.88\linewidth}{\centering ACM Figure Placeholder\\[0.4em]\normalsize Pipeline overview, rendering graph, or system diagram}}
\caption{Representative figure block used to inspect the official ACM manuscript layout, caption spacing, and the balance between body text and the first major figure.}
\end{figure}
\section{Experiments}
The experiments section provides enough structure to expose table treatment, subsection rhythm, and the transition from methods to quantitative comparison.
\subsection{Experimental Setup}
This paragraph acts as a placeholder for datasets, user studies, rendering conditions, baseline systems, and protocol details.
\subsection{Quantitative Results}
\begin{table}[t]
\caption{Quantitative comparison block used to inspect the official ACM manuscript template under realistic submission-style content.}
\centering
\begin{tabular}{lcc}
\toprule
Method & Score & Time \\
\midrule
Placeholder A & 0.91 & 1.4 \\
Placeholder B & 0.88 & 1.7 \\
Placeholder C & 0.84 & 2.0 \\
\bottomrule
\end{tabular}
\end{table}
\section{Conclusion and Discussion}
This page verifies the official ACM template path while preserving the six-part manuscript skeleton requested for the library. The resulting preview is intended to feel closer to a serious review-format submission than to a minimal placeholder sheet.
\begin{thebibliography}{9}
\bibitem{ref1} A. Smith and B. Lee. 2025. Template-aware ACM manuscript preparation. \emph{ACM Journal Placeholder} 1, 1 (2025), 1--10.
\bibitem{ref2} C. Garcia and D. Patel. 2026. Figure and table balance in review-format ACM articles. \emph{ACM Comput. Surv.} 59, 2 (2026), 1--14.
\end{thebibliography}
\end{document}
